% Options for packages loaded elsewhere
\PassOptionsToPackage{unicode}{hyperref}
\PassOptionsToPackage{hyphens}{url}
\PassOptionsToPackage{dvipsnames,svgnames*,x11names*}{xcolor}
%
\documentclass[12pt,
  landscape]{article}
\usepackage{lmodern}
\usepackage{amsmath}
\usepackage{ifxetex,ifluatex}
\ifnum 0\ifxetex 1\fi\ifluatex 1\fi=0 % if pdftex
  \usepackage[T1]{fontenc}
  \usepackage[utf8]{inputenc}
  \usepackage{textcomp} % provide euro and other symbols
  \usepackage{amssymb}
\else % if luatex or xetex
  \usepackage{unicode-math}
  \defaultfontfeatures{Scale=MatchLowercase}
  \defaultfontfeatures[\rmfamily]{Ligatures=TeX,Scale=1}
\setmainfont{IBM Plex Sans}
\setsansfont{Patua One}
\fi
% Use upquote if available, for straight quotes in verbatim environments
\IfFileExists{upquote.sty}{\usepackage{upquote}}{}
\IfFileExists{microtype.sty}{% use microtype if available
  \usepackage[]{microtype}
  \UseMicrotypeSet[protrusion]{basicmath} % disable protrusion for tt fonts
}{}
\makeatletter
\@ifundefined{KOMAClassName}{% if non-KOMA class
  \IfFileExists{parskip.sty}{%
    \usepackage{parskip}
  }{% else
    \setlength{\parindent}{0pt}
    \setlength{\parskip}{6pt plus 2pt minus 1pt}}
}{% if KOMA class
  \KOMAoptions{parskip=half}}
\makeatother
\usepackage[dvipsnames]{xcolor}
\IfFileExists{xurl.sty}{\usepackage{xurl}}{} % add URL line breaks if available
\IfFileExists{bookmark.sty}{\usepackage{bookmark}}{\usepackage{hyperref}}
\hypersetup{
  pdftitle={Reading Group},
  pdfauthor={Scott Cunningham},
  colorlinks=true,
  linkcolor=blue,
  filecolor=Maroon,
  citecolor=Blue,
  urlcolor=Blue,
  pdfcreator={LaTeX via pandoc}}
\urlstyle{same} % disable monospaced font for URLs
\usepackage[a4paper, margin=1in]{geometry}

\usepackage{longtable,booktabs}
\usepackage{calc} % for calculating minipage widths
% Correct order of tables after \paragraph or \subparagraph
\usepackage{etoolbox}
\makeatletter
\patchcmd\longtable{\par}{\if@noskipsec\mbox{}\fi\par}{}{}
\makeatother
% Allow footnotes in longtable head/foot
\IfFileExists{footnotehyper.sty}{\usepackage{footnotehyper}}{\usepackage{footnote}}
\makesavenoteenv{longtable}
\usepackage{graphicx}
\makeatletter
\def\maxwidth{\ifdim\Gin@nat@width>\linewidth\linewidth\else\Gin@nat@width\fi}
\def\maxheight{\ifdim\Gin@nat@height>\textheight\textheight\else\Gin@nat@height\fi}
\makeatother
% Scale images if necessary, so that they will not overflow the page
% margins by default, and it is still possible to overwrite the defaults
% using explicit options in \includegraphics[width, height, ...]{}
\setkeys{Gin}{width=\maxwidth,height=\maxheight,keepaspectratio}
% Set default figure placement to htbp
\makeatletter
\def\fps@figure{htbp}
\makeatother
\setlength{\emergencystretch}{3em} % prevent overfull lines
\providecommand{\tightlist}{%
  \setlength{\itemsep}{0pt}\setlength{\parskip}{0pt}}
\setcounter{secnumdepth}{-\maxdimen} % remove section numbering
% Make \paragraph and \subparagraph free-standing
\ifx\paragraph\undefined\else
  \let\oldparagraph\paragraph
  \renewcommand{\paragraph}[1]{\oldparagraph{#1}\mbox{}}
\fi
\ifx\subparagraph\undefined\else
  \let\oldsubparagraph\subparagraph
  \renewcommand{\subparagraph}[1]{\oldsubparagraph{#1}\mbox{}}
\fi
\pagestyle{empty}
\makeatletter
\@ifpackageloaded{caption}{}{\usepackage{caption}}
\AtBeginDocument{%
\ifdefined\contentsname
  \renewcommand*\contentsname{Table of contents}
\else
  \newcommand\contentsname{Table of contents}
\fi
\ifdefined\listfigurename
  \renewcommand*\listfigurename{List of Figures}
\else
  \newcommand\listfigurename{List of Figures}
\fi
\ifdefined\listtablename
  \renewcommand*\listtablename{List of Tables}
\else
  \newcommand\listtablename{List of Tables}
\fi
\ifdefined\figurename
  \renewcommand*\figurename{Figure}
\else
  \newcommand\figurename{Figure}
\fi
\ifdefined\tablename
  \renewcommand*\tablename{Table}
\else
  \newcommand\tablename{Table}
\fi
}
\@ifpackageloaded{float}{}{\usepackage{float}}
\floatstyle{ruled}
\@ifundefined{c@chapter}{\newfloat{codelisting}{h}{lop}}{\newfloat{codelisting}{h}{lop}[chapter]}
\floatname{codelisting}{Listing}
\newcommand*\listoflistings{\listof{codelisting}{List of Listings}}
\makeatother
\makeatletter
\makeatother
\makeatletter
\@ifpackageloaded{caption}{}{\usepackage{caption}}
\@ifpackageloaded{subcaption}{}{\usepackage{subcaption}}
\makeatother
\ifluatex
  \usepackage{selnolig}  % disable illegal ligatures
\fi

\title{\textbf{Reading Group}}
\usepackage{etoolbox}
\makeatletter
\providecommand{\subtitle}[1]{% add subtitle to \maketitle
  \apptocmd{\@title}{\par {\large #1 \par}}{}{}
}
\makeatother
\subtitle{Lalonde - 40 Years Later: Imbens and Xu's Review}
\author{Scott Cunningham}
\date{2024-07-16}

% my preamble additions =============================================================

% left justify text with ragged right edge =====================================
\usepackage[document]{ragged2e}

% page numbering ===============================================================
\usepackage{fancyhdr}
\pagestyle{fancy}
\fancyhf{}
\renewcommand{\headrulewidth}{0pt}
\fancyfoot[C]{\footnotesize\textcolor{gray}{Reading Group}}
\fancyfoot[R]{\footnotesize\textcolor{gray}{\thepage}}

% change default spacing before each new paragraph (sets as text height) =======
\setlength\parskip{1em plus 0.1em minus 0.2em}

% set formatting for section and subsection headings ===========================
% adding a period after section numbering
% \usepackage{titlesec}
% \titlelabel{\thetitle.\quad}
% comment/uncomment the next line to insert a page break before each new H1
\let\oldsection\section
\renewcommand\section{\clearpage\oldsection}

% changing options for table/figure captions ===================================
% uncomment the below lines to just number figures/tables sequentially
% throughout entire document
% \renewcommand{\thetable}{\thesection.\arabic{table}}
% \renewcommand{\thefigure}{\thesection.\arabic{figure}}

\usepackage[labelfont=it,textfont=it,singlelinecheck=off,justification=raggedright,labelsep=period]{caption}

% resets table/figure counters within each section/subsection ==================
% \usepackage{chngcntr}
% \counterwithin*{table}{section}
%\counterwithin*{table}{subsection}
% \counterwithin*{figure}{section}
%\counterwithin*{figure}{subsection}

\usepackage[breakable]{tcolorbox}

\definecolor{cosmiclatte}{RGB}{255, 248, 231}
\definecolor{justoffwhite}{RGB}{255, 254, 250}

%\definecolor{notes-bg}{RGB}{254, 250, 233}
%\definecolor{notes-frame}{RGB}{222, 217, 193}

%\definecolor{Questions-bg}{RGB}{236, 239, 244}
%\definecolor{Questions-frame}{RGB}{236, 239, 244}

\definecolor{notes-bg}{RGB}{244, 247, 252}
\definecolor{notes-frame}{RGB}{233, 235, 238}
\definecolor{notes-text}{RGB}{47, 54, 61}

\definecolor{slide-titles}{RGB}{220, 50, 47}

\usepackage{pagecolor}
\usepackage{afterpage}

\renewcommand{\baselinestretch}{1.5}

% resize headings =============================

\makeatletter

\renewcommand{\section}{\@startsection
{section}%                   % the name
{1}%                         % the level
{\z@}%                       % the indent / 0mm
{3\baselineskip}%            % the before skip / -3.5ex \@plus -1ex \@minus -.2ex
{3\baselineskip}%          % the after skip / 2.3ex \@plus .2ex
{\sffamily\huge\bfseries}} % the style

\renewcommand{\subsection}{\@startsection
{subsection}%                   % the name
{2}%                         % the level
{\z@}%                       % the indent / 0mm
{3\baselineskip}%            % the before skip / -3.5ex \@plus -1ex \@minus -.2ex
{.75\baselineskip}%          % the after skip / 2.3ex \@plus .2ex
{\sffamily\Large\bfseries\textcolor{slide-titles}}} % the style

\renewcommand{\subsubsection}{\@startsection
{subsubsection}%                   % the name
{3}%                         % the level
{\z@}%                       % the indent / 0mm
{3\baselineskip}%            % the before skip / -3.5ex \@plus -1ex \@minus -.2ex
{.75\baselineskip}%          % the after skip / 2.3ex \@plus .2ex
{\large\bfseries}} % the style

\makeatother

\usepackage[export]{adjustbox}

% ===================================================================================

\begin{document}


      % \maketitle
    \begin{center}
    \vspace*{5cm}
    {\Huge\bfseries Reading Group}\\
        \end{center}  
    % turn off page numbering for titlepage
    \thispagestyle{empty}
        \newpage
  





\newpage

\subsection{Introduction}\label{introduction}

\begin{itemize}
\tightlist
\item
  Experience:

  \begin{itemize}
  \tightlist
  \item
    Initial belief: Solid study = publication.
  \item
    Emphasis on methodology over theory led to rejections.
  \end{itemize}
\item
  Realization:

  \begin{itemize}
  \tightlist
  \item
    Importance of theoretical implications and not just methodological
    contributions.
  \end{itemize}
\item
  Common Mistake:

  \begin{itemize}
  \tightlist
  \item
    Misunderstanding what journals value, leading to repeated
    rejections.
  \end{itemize}
\item
  How to improve submissions?
\end{itemize}

\vfill

\begin{tcolorbox}[beforeafter skip=1cm, ignore nobreak=true, breakable, colframe=notes-frame, colback=notes-bg, coltext=notes-text, boxsep=2mm, arc=0mm, boxrule=0.5mm]

I have already had this experience myself. A few months ago, my research
group submitted a paper to a highly respected conference. In the
double-blind review process, various people criticized us for not having
made a sufficient theoretical contribution.

\end{tcolorbox}

\newpage

\newpage

\subsection{Understanding Journal
Expectations}\label{understanding-journal-expectations}

Understanding the \emph{Publishing Environment}:

\begin{itemize}
\tightlist
\item
  Journal Space is Limited:

  \begin{itemize}
  \tightlist
  \item
    Not all methodologically articles are published.
  \end{itemize}
\item
  Using Strategic Submission:

  \begin{itemize}
  \tightlist
  \item
    Viewing manuscript submission as a strategic game.
  \item
    Aligning manuscripts like firms competing for limited
    resources---journal pages.
  \end{itemize}
\end{itemize}

\vfill

\begin{tcolorbox}[beforeafter skip=1cm, ignore nobreak=true, breakable, colframe=notes-frame, colback=notes-bg, coltext=notes-text, boxsep=2mm, arc=0mm, boxrule=0.5mm]

These points seem obvious at first, but should always be considered.

\end{tcolorbox}

\newpage

\newpage

\subsection{Aligning Manuscript with
Journal}\label{aligning-manuscript-with-journal}

\begin{itemize}
\tightlist
\item
  Two Critical Assumptions:

  \begin{enumerate}
  \def\labelenumi{\arabic{enumi}.}
  \tightlist
  \item
    Manuscript must match its journal's focus.
  \item
    Closer the match, higher the publication likelihood.
  \end{enumerate}
\item
  Key Questions:

  \begin{itemize}
  \tightlist
  \item
    Does your study align with the journal's objectives?
  \item
    How does your manuscript contribute to the theory?
  \end{itemize}
\end{itemize}

\vfill

\begin{tcolorbox}[beforeafter skip=1cm, ignore nobreak=true, breakable, colframe=notes-frame, colback=notes-bg, coltext=notes-text, boxsep=2mm, arc=0mm, boxrule=0.5mm]

The author builds his argumentation mainly on these two points, whereby
the two key questions for his own work arise. As for the second
question, I am personally unsure whether it always makes sense, but more
on that in the discussion.

\end{tcolorbox}

\newpage

\newpage

\subsection{Evaluating the
Contribution}\label{evaluating-the-contribution}

Three Strategic Tests for Contributions:

\begin{enumerate}
\def\labelenumi{\arabic{enumi}.}
\tightlist
\item
  Value: Does the contribution add meaningful insights?
\item
  Imitability: Is the contribution unique to the theory being advanced?

  \begin{itemize}
  \tightlist
  \item
    Contribution differ from existing theories?
  \end{itemize}
\item
  Rareness: Does the contribution provide a novel insight?

  \begin{itemize}
  \tightlist
  \item
    Originality?
  \end{itemize}
\end{enumerate}

\vfill

\begin{tcolorbox}[beforeafter skip=1cm, ignore nobreak=true, breakable, colframe=notes-frame, colback=notes-bg, coltext=notes-text, boxsep=2mm, arc=0mm, boxrule=0.5mm]

The author recommends applying the following three tests to your own
work. What about a paper that implements theoretical constructs as
methods so that they can be applied? For example, many machine learning
variants arise from their implementation and the associated
possibilities (LightGBM, XGBoost, ExtraTrees)

\end{tcolorbox}

\newpage

\newpage

\subsection{Discussion Starters}\label{discussion-starters}

Lets talk about the optinion from Bergh.

\begin{itemize}
\tightlist
\item
  Do Bergh's points apply to all areas of research?

  \begin{itemize}
  \tightlist
  \item
    Some areas are very fast-moving (e.g.~research in ML and AI) and
    more applied in general.
  \end{itemize}
\item
  Despite possible tests: The evaluation of whether a theoretical
  contribution is relevant and a real novelty is very subjective and
  depends on both the opinion and the specific research area of the
  reviewer.

  \begin{itemize}
  \tightlist
  \item
    For some general topics, this may be less relevant, but I have
    sometimes experienced a large variance in the judgment of reviewers
    from different disciplines when it comes to specific contributions
    at scientific conferences.
  \end{itemize}
\end{itemize}

\newpage

\newpage

\vspace*{2.5cm}

\begin{center}

\section{Thanks for the attention!}\label{thanks-for-the-attention}

\end{center}

\end{document}
